\documentclass[fc]{tarea}

\usepackage[style=mexican]{csquotes}

\newtheorem{definition}{Definición}[chapter]
\newtheorem{theorem}{Teorema}[chapter]

\newcommand{\al}{\symscr{A}}
\newcommand{\lan}{\symscr{L}}
\newcommand{\sis}{\symup{SF}}
\newcommand{\expresiones}{\symup{Exp}}
\NewDocumentCommand{\expr}{O{\al}}{\expresiones(#1)}

%conjuntos
\providecommand\st{}
\newcommand\SetSymbol[1][]{%
\nonscript\:#1\vert
\allowbreak
\nonscript\:
\mathopen{}}
\DeclarePairedDelimiterX\set[1]{\{}{\}}{%
\renewcommand\st{\SetSymbol[\delimsize]}
#1
}

\begin{document}

\chapter*{Estructuras}
\begin{definition}
Un tipo de semejanza $\rho$ está formado por un conjunto de constantes $\mathcal{C}$, un conjunto de letras funcionales $\mathcal{F}_n$ de aridad $n$, para cada 
$n \in \mathbb{N}$ y un conjunto de letras predicativas $\mathcal{R}_n$ de aridad $n$, para cada $n \in \mathbb{N}$. Es decir,

$$\rho = \mathcal{C} \cup \bigcup_{n \in \mathbb{N}} \mathcal{F}_n \cup \bigcup_{n \in \mathbb{N}} \mathcal{R}_n$$
\end{definition}

El tipo de semejanza son todos aquellos símbolos que le ``dan estructura'' a un conjunto. Pero necesitamos hablar de lo que pueden hacer, es decir, el tipo de
semejanza nos da la posibilidad y el lengaje nos permite describirlos. Con un lenguaje podemos hablar de elementos y evaluar si son verdaderas o falsas.

\begin{definition}
Un lenguaje, de primer orden, $\mathcal{L}$ consta de lo siguiente:

\begin{enumerate}
    \item Un tipo de semejanza $\rho$.
    \item Un conjunto (numerable) de variables llamado $VAR$.
    \item Símbolos lógicos $\{\neg, \land, \lor, \to, \leftrightarrow, \forall, \exists\}$.
    \item Símbolos auxiliares $\{(, ), ', \dots\}$.
    \item Igualdad $\{=\}$.
\end{enumerate}
\end{definition}

Todos los lenguajes coinciden en los puntos 2–5, entonces dos lenguajes son diferentes si difieren en su tipo de semejanza. Por esta razón muchos autores 
prefieren llamar lenguaje al tipo de semejanza y obviar que tiene todos los demás símbolos.

\begin{definition}
  Un término consiste en lo siguiente:
  \begin{enumerate}
    \item Las variables y constantes son términos.
    \item Si $t_1, t_2, \dots, t_n$ son términos y $f\in \mathcal{F}_n$, entonces $f(t_1,\cdots, t_n)$ es un término.
  \end{enumerate}
\end{definition}

\begin{definition}
  Sea $A$ un conjunto distinto del vacío, diremos que $A$ es el universo de interpretación. Una función de interpretación, en $A$, es una función 
  $I:\rho \to A \cup \bigcup_{n \in \mathbb{N}} \mathcal{P}(A^n)$ que cumple lo siguiente:
  \begin{enumerate}
    \item Si $c \in \mathcal{C}$, entonces $I(c)\in A$.
    \item Si $f \in \mathcal{F}_n$, entonces $I(f): A^n \to A$.
    \item Si $R \in \mathcal{R}_n$, entonces $I(R)\subset A^n$.
  \end{enumerate}
\end{definition}

\begin{definition}
  Una estructura es una pareja $(A, I)$ donde $A$ es un conjunto no vacío e $I$ es una función de interpretación en $A$. Esta pareja se denota $\mathfrak{A} = (A,I)$.
\end{definition}

\end{document}